\section{Introduction}
\label{sec:intro}

To serve resource-constrained heterogeneous multi-task query streams, a prominent paradigm is parameter-efficient fine-tuning via task-specific expert adapters (e.g., LoRA) on a shared frozen base model \cite{hu2021lora, dettmers2023qlora}. At test-time, the system dynamically blends adapter outputs on-the-fly to handle sequential queries sample-by-sample and layer-by-layer.

Stateless dynamic routers like SABLE \cite{sable} and SPS-ZCA \cite{spszca} calculate expert weights at early layers using cosine similarity. However, they suffer from a severe \emph{routing jitter paradox}: because weights are computed independently for each sample and layer, representation noise causes coefficients to oscillate wildly layer-by-layer, corrupting representation manifolds and degrading robustness.

To smooth out oscillations, stateful ensembling has emerged. ChemMerge \cite{chemmerge} models weights via ODE kinetics, while Momentum-Merge \cite{momentum_merge} simplifies this as open-loop Exponential Moving Average (EMA). While stateful methods suppress jitter under stationary streams, they introduce a critical failure mode: \emph{inertial drag} (phase delay). Accumulating history in an open-loop fashion, they adapt slowly to task switches, collapsing accuracy under heterogeneous (rapidly switching) workloads (e.g., dropping to $88.42\%$ and $86.17\%$). Additionally, continuous-time kinetics require costly ODE solvers, violating latency budgets.

We resolve this stability-responsiveness trade-off by introducing \textbf{PID-Merge}, a closed-loop, discrete-time Proportional-Integral-Derivative (PID) controlled stateful routing framework. Rooted in classical control theory \cite{astrom_pid, ziegler1942optimum}, PID-Merge treats raw early-layer weights as the setpoint and active coefficients as the plant output. By feeding back the tracking error utilizing Proportional (P), Integral (I), and Derivative (D) gains, PID-Merge achieves exceptional smoothing while eliminating tracking lag. Crucially, the \textbf{Derivative (D) term} measures error acceleration to instantly detect task transitions and provide an anticipatory boost driving weights to the target within 2--3 layers. For deep topologies, our control-theoretic \textbf{anti-windup clamping} (conditional integration) freezes error accumulation near boundary saturation, ensuring robust responsiveness.

PID-Merge is highly lightweight and deployment-ready: (1) \textbf{Minimal Serving Latency:} It uses an incremental velocity algorithm running in $O(1)$ time, requiring only $15$ FLOPs per expert-layer. (2) \textbf{Zero-Shot Readiness:} It achieves outstanding results using robust heuristic default gains ($K_p = 0.5, K_i = 0.15, K_d = 0.2$). (3) \textbf{Calibration Efficiency:} On a tiny 32-sample calibration sequence, the gains are optimized via gradient descent with stable out-of-sample generalization.

\paragraph{Our Contributions.}
1. We propose \textbf{PID-Merge}, the first closed-loop, discrete-time PID control-theoretic framework for dynamic stateful model serving, and show how the Derivative (D) term suppresses representation noise while completely eliminating tracking lag. To ensure scalability to deep topologies, we integrate a control-theoretic \textbf{anti-windup clamping mechanism} that prevents saturation-induced transition delays.
2. We evaluate PID-Merge on the Isolating Coordinate Sandbox (ICS). On overlapping heterogeneous streams, PID-Merge achieves up to \textbf{94.82\%} accuracy—yielding up to \textbf{+6.40\%} improvement over SOTA ChemMerge and \textbf{+8.65\%} over Momentum-Merge—while slashing routing jitter and running with negligible latency.
3. We physically validate PID-Merge on a 12-layer GPT-2 Small backbone routing 3 actual fine-tuned task adapters on an NVIDIA A100 GPU (Appendix~\ref{sec:appendix_gpt2}). The physical results confirm that PID-Merge slashes depth-wise jitter by over \textbf{73\%} while maintaining near-oracle serving accuracy, adding an imperceptible latency overhead of just \textbf{0.012 ms} ($40\times$ faster than SOTA ChemMerge).
